\section{Colocated vs. Prefill or Decode Composition}

This section documents how the runtime composes the stage-cost formulas into
request-visible TTFT, TPOT, E2E, and goodput. The relevant implementations live
in \texttt{astra-sim/workload/ColocatedServingRuntime.cc},
\texttt{astra-sim/workload/PdServingRuntime.cc}, and
\texttt{astra-sim/workload/serving/topology/ServingTopology.cc}.

\subsection{Replica Assignment and Admission Domains}

Let \(R\) be the configured \texttt{dp\_replica\_count}. The current replica
mapper is:

\begin{equation*}
\mathrm{replica}(r) =
\left(
\mathrm{request\_id}(r) + \mathrm{original\_index}(r)
\right)
\bmod R
\end{equation*}

This is the exact policy implemented by
\texttt{ServingTopology::assign\_replica(...)}.

For colocated serving, the per-replica admission budget is:

\begin{equation*}
M_{\mathrm{replica}} =
\left\lceil
\frac{M_{\mathrm{global}}}{R}
\right\rceil
\end{equation*}

where \(M_{\mathrm{global}}\) is
\texttt{scheduler.max\_running\_requests}. This is the exact logic used by
\texttt{ColocatedServingRuntime::max\_running\_requests\_per\_replica()}.

\subsection{Colocated Request Composition}

In the colocated runtime, a request first waits for prefill admission, then
executes prefill, then may wait again for its first decode batch. The exact
request-level TTFT composition implied by the recorded timestamps is:

\begin{equation*}
\mathrm{TTFT}_{\mathrm{colocated}} =
W_p + T_{\mathrm{prefill}} + W_d + T_{\mathrm{decode}}^{(1)}
\end{equation*}

The compact architecture summary used in planning discussions is:

\begin{equation*}
\mathrm{TTFT}_{\mathrm{colocated}} \approx
W_q + T_{\mathrm{prefill}} + T_{\mathrm{decode}}^{(1)}
\end{equation*}

The exact end-to-end composition is:

\begin{equation*}
\mathrm{E2E}_{\mathrm{colocated}} =
W_p + T_{\mathrm{prefill}} + W_d + T_{\mathrm{decode}}^{(\mathrm{all})}
\end{equation*}

\subsection{Prefill or Decode Disaggregation}

In the PD runtime, the request may wait in distinct prefill, transfer, and
decode queues. The exact TTFT composition is:

\begin{equation*}
\mathrm{TTFT}_{\mathrm{PD}} =
W_p + T_{\mathrm{prefill}} +
W_t + T_{\mathrm{transfer}} +
W_d + T_{\mathrm{decode}}^{(1)}
\end{equation*}

The exact end-to-end composition is:

\begin{equation*}
\mathrm{E2E}_{\mathrm{PD}} =
W_p + T_{\mathrm{prefill}} +
W_t + T_{\mathrm{transfer}} +
W_d + T_{\mathrm{decode}}^{(\mathrm{all})}
\end{equation*}

\subsection{KV Placement Under DP-Attention}

After prefill or transfer completes, the runtime stores a shard-local KV size:

\begin{equation*}
Q_{\mathrm{KV,resident}} =
\frac{Q_{\mathrm{KV}}}{d_a}
\end{equation*}

where \(d_a\) is the decode-side \texttt{dp\_attention\_degree}. This is the
exact rule used in both colocated and PD runtimes when recording
\texttt{kv\_resident\_bytes}.

\subsection{How Topology Changes Goodput}

The topology features change goodput through two paths:

\begin{equation*}
\mathrm{TTFT} \uparrow \text{ or } \downarrow
\Longleftarrow
\left\{
\begin{array}{l}
\text{replica-local admission}, \\
\text{batch shape}, \\
\text{TP, PP, EP, or DP-attention costs}, \\
\text{optional PD transfer}
\end{array}
\right\}
\end{equation*}

\begin{equation*}
\mathrm{TPOT} \uparrow \text{ or } \downarrow
\Longleftarrow
\left\{
\begin{array}{l}
\text{decode batch size}, \\
\text{decode layout}, \\
\text{decode-side communication terms}, \\
\text{decode queue pressure}
\end{array}
\right\}
\end{equation*}

Since goodput is defined by the SLO predicate in Section~1, any topology change
that alters TTFT, TPOT, or E2E changes the good-request count and therefore the
reported goodput.
